\section{Why the Contract Arm Wins}
\label{sec:why}

Aggregate accuracy does not explain itself. This section traces one task end
to end, reads the SQL every arm submitted, and accounts for the work and
money each arm spent.

\subsection{The rule was in both prompts; only one arm reached it}
\label{sec:reach}

Task~1278 asks for \emph{the average fee NexPay would charge for a credit
transaction of 50 EUR}.

\begin{center}
\footnotesize
\setlength{\tabcolsep}{4pt}
\begin{tabular}{@{}lr@{}}
\toprule
Gold      & \texttt{0.352294} \\
\armC{} (correct)  & \texttt{0.352294} \\
\armM{} (wrong)    & \texttt{0.353053} \\
\armS{} (wrong)    & \texttt{0.353053} \\
\midrule
\texttt{is\_credit IS NULL OR is\_credit = true} & \texttt{0.352294} \\
\texttt{is\_credit = true} & \texttt{0.353053} \\
\bottomrule
\end{tabular}
\end{center}

Both figures reproduce against the warehouse with the two predicates
shown. A \texttt{NULL} in a \texttt{fees} column means \emph{this rule
applies to every value}, not \emph{unknown}; both baselines dropped the
\texttt{NULL} rows and produced the identical wrong number. The contract arm called
\texttt{lookup\_domain("fees")}, and its recorded reasoning states the rule
back: \emph{``rules that apply to credit transactions: is\_credit true or
null''}.

The decisive detail is that \texttt{manual.md} documents this rule too.
\armM{} had it, inside a 24{,}177-character system prompt. \armC{}'s system
prompt is 2{,}475 characters, with the rule reachable by one targeted call.
Two published results frame the difference: retrieval quality falls for
material buried mid-prompt~\cite{lost-in-the-middle}, and what a tool's
documentation says changes what an agent does with
it~\cite{tool-documentation}. \armM{}'s copy of the rule was buried
mid-prompt; \armC{}'s was a tool's documentation.

The finding is not that the contract arm was given more information. For
\emph{this} rule both arms had it, and only one used it. What the contract
changes is not what is knowable but what is reachable at the moment of use.
That scoping matters: it holds for the clauses the manual states outright,
which is the class this failure belongs to, and not for the three clauses
the contract resolves on its own (Section~\ref{sec:threats}).

\subsection{The contract's clauses reach the SQL}
\label{sec:clauses}

The claim that delivery matters as much as possession can be observed
directly, by reading the SQL each agent submitted rather than the answer it
produced. We defined one detector per clause of the contract's fee semantics
that a correct query must express: the \texttt{NULL}-wildcard disjunct, the
empty-list wildcard, the capture-delay band mapping, a per-merchant monthly
aggregate, the natural-month reconstruction, and fraud measured on euro
volume. We then restricted the comparison to the 97 tasks that require all
six (69 on \mDS{}, whose truncation cost the rest), so every row compared
carries the same requirement (Table~\ref{tab:clauses}). The detectors are
deliberately permissive: the question is whether the agent expressed the
idea, not whether it copied the contract's phrasing.

\begin{table}[t]
\centering
\footnotesize
\setlength{\tabcolsep}{3pt}
\begin{tabular}{@{}lrrrr@{}}
\toprule
\tabhead{Wrote all six clauses} & \mGLM{} & \mDS{} & \mSON{} & \mGPT{} \\
\midrule
\armS{} & 7\% & 4\% & 3\% & 8\% \\
\armH{} & 14\% & 4\% & 4\% & 8\% \\
\armM{} & 3\% & 9\% & 4\% & 4\% \\
\armC{} & \textbf{39\%} & \textbf{65\%} & \textbf{55\%} & \textbf{98\%} \\
\bottomrule
\end{tabular}
\caption{Share of the 97 tasks requiring all six fee clauses (69 on \mDS{})
on which the agent's submitted SQL expresses all six. \armM{} holds the same
knowledge as prose in its system prompt. \armC{} against \armM{} is
$p{=}2{\times}10^{-10}$, $3{\times}10^{-12}$, $9{\times}10^{-16}$ and
$2{\times}10^{-47}$ (Fisher exact).}
\label{tab:clauses}
\end{table}

\textbf{Note what does not move.} \armM{} writes no more of the contract's
clauses on a strong model than on a weak one (3\%, 9\%, 4\%, 4\%), while
\armC{} reaches 55\% and 98\% on the two frontier models. \mSON{} is the
extreme case: 4\% against 55\% in the SQL, no better than a bare schema
against 45 points above it in accuracy. Capability does not, on its own,
make a model extract structure from prose; it makes a model much better at
\emph{using} structure it is handed. That is the same asymmetry the
derivation gap shows (Section~\ref{sec:ceiling}), measured on behaviour
instead of outcomes.

\textbf{What this does not measure.} Clause presence does not predict
success: within every arm on every model, attempts that got the task right
and attempts that got it wrong write the same clauses (every Fisher
$p{\ge}0.14$), and \armH{} on \mGLM{} writes more clauses than \armM{} while
scoring at the bare-schema floor. The measure captures whether the
contract's vocabulary reaches the query, not whether the query is any good.

\subsection{Less work, and usually less money}
\label{sec:efficiency}

\begin{table*}[t]
\centering
\small
\begin{tabular}{@{}lrrrr@{}}
\toprule
 & \armS{} & \armH{} & \armM{} & \armC{} \\
\midrule
\multicolumn{5}{@{}l}{\emph{\mGLM{}}}\\
Turns / task & 12.1 & 8.4 & 9.8 & \textbf{7.1} \\
Tool calls & 5{,}134 & 4{,}110 & 4{,}282 & \textbf{3{,}355} \\
Reasoning tokens & 428k & 525k & 400k & \textbf{154k} \\
Cost / correct & \$0.0084 & \$0.0080 & \$0.0076 & \textbf{\$0.0029} \\
Forced answers & 41 & 0 & 18 & \textbf{0} \\
\midrule
\multicolumn{5}{@{}l}{\emph{\mDS{}}}\\
Turns / task & 13.7 & 11.4 & 12.8 & \textbf{8.2} \\
Tool calls & 5{,}489 & 6{,}375 & 4{,}770 & \textbf{4{,}155} \\
Reasoning / task & 10{,}941 & 10{,}766 & 9{,}243 & \textbf{4{,}844} \\
Cost / correct & \$0.0195 & \$0.0173 & \$0.0117 & \textbf{\$0.0061} \\
Forced answers & 49 & 15 & 38 & \textbf{0} \\
\midrule
\multicolumn{5}{@{}l}{\emph{\mSON{}}}\\
Turns / task & 11.3 & 10.0 & 11.1 & \textbf{7.0} \\
Tool calls & 5{,}644 & 6{,}494 & 5{,}176 & \textbf{3{,}791} \\
Reasoning tokens & 942k & 1{,}411k & 1{,}044k & \textbf{824k} \\
Cost / correct & \$0.288 & \$0.259 & \$0.311 & \textbf{\$0.136} \\
Forced answers & 43 & 27 & 32 & \textbf{0} \\
\midrule
\multicolumn{5}{@{}l}{\emph{\mGPT{}}}\\
Turns / task & 7.9 & 10.4 & 7.6 & \textbf{6.5} \\
Tool calls & 5{,}099 & 7{,}997 & \textbf{4{,}354} & 5{,}245 \\
Reasoning tokens & 357k & 311k & 277k & \textbf{141k} \\
Cost / correct & \$0.082 & \$0.099 & \textbf{\$0.062} & \$0.067 \\
Forced answers & 0 & 8 & 0 & \textbf{0} \\
\bottomrule
\end{tabular}
\caption{Work and cost per arm. The contract arm uses the fewest turns and
the least reasoning on every model and never exhausts its output budget; it
is the cheapest arm per correct answer on three of the four. \mDS{}'s
reasoning row is per task rather than a total, because its complete-case set
has fewer rows than the other runs.}
\label{tab:efficiency}
\end{table*}

\armC{} is not thinking harder; it has less to search for. On all four
models it uses the fewest turns and the least reasoning (by a factor of two
to three on three models, by 12\% on \mSON{}), and it is the only arm that
never exhausted its output budget on any model (Table~\ref{tab:efficiency};
on reporting cost beside accuracy, see~\cite{agents-that-matter}).

Money follows work on three models and not on the fourth. \armC{} is the
cheapest arm outright on both flash models, and per correct answer it is the
cheapest on \mGLM{}, \mDS{} and \mSON{}, by about $2{\times}$ on \mSON{}
(\$0.136 against \$0.259 for the next arm). On \mGPT{} it bills \$21.32
against \armM{}'s \$14.42, and \$0.067 per correct answer against \$0.062.
The recorded token counts locate the difference (Table~\ref{tab:gptcost}):
fresh, uncached input accounts for \$5.3 of the \$6.9 gap and output for
\$1.6, and cached input is a wash. \armC{} reads 1{,}692 fresh tokens per
turn against \armM{}'s 582, because its knowledge arrives as tool results.
On \mGPT{} it calls \texttt{lookup\_metric} 1{,}690 times and
\texttt{lookup\_domain} 776 times over 401 tasks, two tool calls per turn,
and each result is billed at the fresh rate when it first enters the
context. \armM{}'s knowledge sits in a system prompt that is the same prefix
in every request of every task, and is billed at the cached rate. \armC{}
also emits more visible output (726k tokens against 425k) despite half the
reasoning (141k against 277k), because \mGPT{} validates before it runs: 676
\texttt{inspect\_query} calls against 622 \texttt{run\_query}, so most
queries are written twice. \mSON{} shows the same mechanism (1{,}976 against
1{,}037 fresh tokens per turn), but its \armM{} arm carries 191k tokens of
context per task and pays \$14.40 in cached input alone, which is what makes
the contract arm cheaper there.

\begin{table}[t]
\centering
\footnotesize
\setlength{\tabcolsep}{5pt}
\begin{tabular}{@{}lrrrr@{}}
\toprule
\mGPT{} & \tabhead{Fresh input} & \tabhead{Cached input} & \tabhead{Output} & \tabhead{Total} \\
\midrule
\armM{} & \$3.55 & \$3.85 & \$7.02 & \$14.42 \\
\armC{} & \$8.87 & \$3.78 & \$8.67 & \$21.32 \\
\bottomrule
\end{tabular}
\caption{\mGPT{}'s bill for the two content-bearing arms, decomposed from the
recorded token counts at the pinned endpoint's prices: \$2.00 per million
fresh input tokens, \$0.20 cached, \$10.00 output (reasoning included).}
\label{tab:gptcost}
\end{table}

The general statement is about prices, not about the contract: at a
ten-to-one ratio between fresh and cached input, on-demand delivery is
charged per task and a static prompt once per cache window, and whether the
contract's savings in turns and reasoning cover the difference depends on
how often the model looks things up and how long its conversations run.
What generalises across the four models is the work claim, not the price
claim.

\subsection{Where the contract arm fails}
\label{sec:losses}

On \mGLM{} \armC{} loses 14 tasks to \armS{} and 29 to \armM{}, 10 to both,
and two modes account for nearly all of them: set membership under
\texttt{NULL} wildcards, the rule of Section~\ref{sec:reach} failing in the
other direction (on task~1500, \emph{fee IDs that apply to account\_type = O
and aci = C}, \armC{} returns both false positives and false negatives while
\armM{} matches exactly), and near-miss arithmetic (task~1274: gold
\texttt{0.126459}, \armC{} \texttt{0.125516}). These are not retrieval
failures. \texttt{lookup\_domain} is called on 331 of 401 tasks, the 70 that
skip it are easier for every arm, and 26 of the 33 losses occur on tasks
where the lookup was performed; \mSON{} repeats the pattern: of the 42 tasks it
loses to at least one other arm, 35 had the lookup. The agent fetches the rule and misapplies it.
That points away from hoisting rule bodies into the system prompt, since
the rules were already fetched, and toward making the wildcard semantics
executable: a
five-line DuckDB macro encoding the rule the contract states in prose
answers tasks~1500 and~1502 exactly. Section~\ref{sec:conclusion} says why
that tool was not built before this measurement.
